% ============================================================
%  Section 4: Building a Data Warehouse (ETL, Data Quality, Data Wrangling)
% ============================================================
\section{Building a Data Warehouse}

% -------- subsection: Process Overview --------

\begin{frame}
  \frametitle{How a Data Warehouse is Built}

  \begin{block}{The Process: Iterative by Nature}
    Building a DWH is \textbf{never a one-time project} but an ongoing,
    iterative effort.
    New requirements, changed source systems, and evolving business needs
    continuously shape the DWH.
  \end{block}

  \begin{block}{The 7 Phases (based on the Rapid Warehousing Methodology)}
    \begin{enumerate}
      \item \textbf{Qualification}: Is the organisation ready for a DWH project?
      \item \textbf{Requirements Analysis}: What data and analyses are needed?
      \item \textbf{Design}: Architecture, data models, process design
      \item \textbf{Implementation}: Build ETL processes, load data, develop marts
      \item \textbf{Acceptance Test}: Validate quality and completeness
      \item \textbf{Deployment}: Train users, go live
      \item \textbf{Review}: Lessons learned, plan next iteration
    \end{enumerate}
  \end{block}
\end{frame}

% -------- subsection: Source Analysis --------

\begin{frame}
  \frametitle{Requirements Analysis: What Do We Need?}

  \begin{block}{Goal}
    Understand \textbf{what questions the business needs to answer}
    and what data is required to answer them.
  \end{block}

  \begin{block}{Key Activities}
    \begin{itemize}
      \item Interview stakeholders: managers, analysts, power users
      \item Identify \textbf{KPIs} and reporting requirements
      \item Determine required \textbf{grain} (day, week, individual transaction)
      \item Define the \textbf{history depth} needed (1 year? 10 years?)
      \item Understand the required \textbf{refresh frequency}:
        daily batch, near-real-time, or on-demand
    \end{itemize}
  \end{block}

  \begin{examples}{Good Questions to Ask Stakeholders}
    \begin{itemize}
      \item \textit{``What decisions does your department make every week?''}
      \item \textit{``What data do you currently look up manually?''}
      \item \textit{``What would you do if you had this data instantly?''}
    \end{itemize}
  \end{examples}
\end{frame}

% ------------------------------------------------------------
\begin{frame}
  \frametitle{Source System Analysis}

  \begin{block}{Goal}
    Understand the \textbf{data landscape}: what exists, where it lives,
    how good it is.
  \end{block}

  \begin{block}{DWH Source Systems Can Be Very Diverse}
    \begin{itemize}
      \item Relational databases: Oracle, DB2, SQL Server, PostgreSQL
      \item ERP / enterprise systems: SAP R/3, SAP BW, Siebel
      \item Legacy systems: IMS, VSAM, COBOL files
      \item Web / Internet sources: APIs, web scraping
      \item Flat files: Excel, CSV, Access, Lotus
      \item Unstructured text files
    \end{itemize}
  \end{block}

  \begin{alertblock}{Common Surprise}
    Data quality in source systems is \textbf{almost always worse} than
    anyone expects. Budget significant effort for data cleansing.
  \end{alertblock}
\end{frame}

% -------- subsection: ETL Process --------

\begin{frame}
  \frametitle{The ETL Process}

  \begin{block}{ETL = Extract, Transform, Load}
    ETL is the central \textbf{data pipeline} that moves data from source
    systems into the DWH.
  \end{block}

  % PICTURE PLACEHOLDER
  \begin{alertblock}{Picture: ETL Flow with Staging Area}
    \textit{Add a three-layer flow diagram:
    (Bottom) Operational source systems
    $\xrightarrow{E}$ Staging Area (raw landing zone)
    $\xrightarrow{T}$ Transformation (cleanse, harmonise, enrich)
    $\xrightarrow{L}$ Data Warehouse (storage + data marts).
    Use arrows labeled E, T, L. Add the ``DWH Manager / Monitor'' below
    as a control element.}
  \end{alertblock}

  \begin{block}{Key Principle: Decouple Extraction from Transformation}
    Extracting data into the staging area first, before transforming it,
    allows \textbf{validation, profiling, and error handling} without
    impacting either the source system or the target DWH.
  \end{block}
\end{frame}

% -------- subsection: Extraction --------

\begin{frame}
  \frametitle{Extract: Getting the Data Out}

  \begin{block}{The First Step}
    Extraction creates access to the operational data sources and copies
    data into the staging area.
  \end{block}

  \begin{block}{Extraction Methods}
    \begin{tabularx}{\linewidth}{@{}l X@{}}
      \toprule
      \textbf{Method}             & \textbf{Description} \\
      \midrule
      Full Refresh (Data Refresh) & Load all data from scratch on each run \\
      Delta / Incremental Update  & Load only changes since the last run \\
      \bottomrule
    \end{tabularx}
  \end{block}

  \begin{block}{Change Data Capture (CDC)}
    \begin{itemize}
      \item \textbf{Timestamp columns}: filter rows where \texttt{UPDATED\_AT > last\_run}
      \item \textbf{Database triggers}: source system records changes automatically
      \item \textbf{Transaction log mining}: read the DB log to find changed rows
      \item \textbf{Full diff}: compare snapshots (last resort, expensive)
    \end{itemize}
  \end{block}
\end{frame}

% ------------------------------------------------------------
\begin{frame}
  \frametitle{Extract: Light Transformations During Extraction}

  Some minor transformations can already be applied during extraction.
  These are technical in nature and do not touch business logic.

  \begin{block}{Acceptable During Extraction}
    \begin{itemize}
      \item Character encoding conversion: ASCII to/from EBCDIC
      \item Decompression
      \item Case conversion (usually to upper case)
      \item Date and time format conversion
      \item Sorting
    \end{itemize}
  \end{block}

  \begin{alertblock}{Warning}
    Applying too much transformation during extraction can:
    \begin{itemize}
      \item \textbf{Obscure errors} --- bad data should be visible in staging
      \item \textbf{Impact performance} of the operational source system
    \end{itemize}
    Keep extraction \textit{as light as possible}.
  \end{alertblock}
\end{frame}

% -------- subsection: Data Quality --------

\begin{frame}
  \frametitle{Data Quality: The Biggest Challenge}

  \begin{block}{Why Data Quality Matters}
    \begin{itemize}
      \item \textbf{40\,\%} of companies have suffered losses due to poor data quality
        \hfill\small(The Data Warehousing Institute)
      \item \textbf{15--20\,\%} of values in typical customer databases are incorrect
        \hfill\small(Uni Trier)
      \item \textbf{50\,\%} of CRM projects fail because of data quality problems
        \hfill\small(Wayne Eckerson, TDWI)
    \end{itemize}
  \end{block}

  \begin{block}{Common Data Quality Issues}
    \begin{tabularx}{\linewidth}{@{}l X@{}}
      \toprule
      \textbf{Issue}       & \textbf{Description} \\
      \midrule
      Completeness         & Missing or NULL values \\
      Consistency          & Same entity represented differently across sources \\
      Accuracy             & Values that are factually wrong or outdated \\
      Timeliness           & Data that arrives too late to be useful \\
      Duplicates           & Same record represented multiple times \\
      \bottomrule
    \end{tabularx}
  \end{block}
\end{frame}

% ------------------------------------------------------------
\begin{frame}
  \frametitle{Data Defect Classification}

  \begin{block}{Three Classes of Defects}
    \begin{tabularx}{\linewidth}{@{}l X X X@{}}
      \toprule
      & \textbf{Class 1}            & \textbf{Class 2}                  & \textbf{Class 3} \\
      \midrule
      Detection  & Automatic               & Automatic                         & Manual \\
      Correction & Automatic               & Manual                            & Manual \\
      \midrule
      Syntactic  & Known format errors     & Detectable format incompatibilities & --- \\
      Semantic   & Missing values (fill)   & Outliers / implausible values       & Undetected semantic errors in source \\
      \bottomrule
    \end{tabularx}
  \end{block}

  \begin{alertblock}{Class 3 Defects}
    Class 3 defects (undetected semantic errors) \textbf{must be fixed at the source
    system} --- they cannot be reliably corrected in the ETL pipeline.
    This requires cooperation with the source system owners.
  \end{alertblock}
\end{frame}

% -------- subsection: Data Wrangling / Transformation --------

\begin{frame}
  \frametitle{Data Wrangling: What is it?}

  \begin{block}{Definition}
    \textbf{Data Wrangling} (also called data munging) is the process of
    transforming and mapping raw data into a format suitable for analysis.
    In the DWH context, this is the heart of the \textbf{Transform} step in ETL.
  \end{block}

  \begin{block}{Typical Wrangling Activities}
    \begin{itemize}
      \item \textbf{Filtering}: Remove irrelevant or erroneous rows
      \item \textbf{Standardisation}: Unify naming conventions and formats
      \item \textbf{Harmonisation}: Align data from multiple sources into one consistent schema
      \item \textbf{Enrichment}: Add derived or calculated attributes
      \item \textbf{Aggregation}: Pre-compute summary values
      \item \textbf{Deduplication}: Identify and remove duplicate records
    \end{itemize}
  \end{block}
\end{frame}

% ------------------------------------------------------------
\begin{frame}
  \frametitle{Data Standardisation: Why it Matters}

  \begin{block}{The Problem}
    The same company can appear in data under many different names:
  \end{block}

  \begin{examples}{Revenue Report Without Standardisation}
    \begin{tabularx}{\linewidth}{@{}l r l@{}}
      \toprule
      \textbf{Company Name}         & \textbf{Revenue (\euro{})} & \textbf{Invoice \#} \\
      \midrule
      BMW                           & 15,687       & 123-65 \\
      Bayrische Motorenwerke        & 47,657       & 875-5  \\
      BMW                           & 134,982      & 123-410 \\
      b.m.w.                        & 9,807        & 459-60 \\
      Bayerische Motoren Werke      & 237,897      & 164-45 \\
      \midrule
      \texttt{WHERE Company = 'BMW'} & \textbf{150,669} & (only 2 rows!) \\
      After standardisation         & \textbf{446,030} & (all 5 rows) \\
      \bottomrule
    \end{tabularx}
  \end{examples}
\end{frame}

% ------------------------------------------------------------
\begin{frame}
  \frametitle{Filtering and Harmonisation}

  \begin{block}{Filtering in the Staging Area}
    After extraction, validation routines are run to detect and correct
    defects automatically where possible.
    Both syntactic and semantic defects can be found.
  \end{block}

  \begin{block}{Harmonisation}
    The process of \textbf{business-level reconciliation} of filtered data.
    After technical standardisation, values from different sources are
    aligned to the same business meaning:
    \begin{itemize}
      \item ``M'' and ``male'' and ``männlich'' $\to$ \texttt{M}
      \item ``DE'', ``Germany'', ``Deutschland'', ``DEU'' $\to$ \texttt{DE}
      \item Date formats: \texttt{15/03/2024}, \texttt{2024-03-15},
        \texttt{20240315} $\to$ ISO 8601: \texttt{2024-03-15}
    \end{itemize}
  \end{block}
\end{frame}

% ------------------------------------------------------------
\begin{frame}
  \frametitle{Deduplication: Matching Data Records}

  \begin{block}{The Problem}
    The same real-world entity appears multiple times in the data ---
    with slight variations in spelling, format, or completeness.
  \end{block}

  \begin{block}{Duplicate Matching Steps}
    \begin{enumerate}
      \item Define \textbf{match criteria}: which fields to compare?
        (name, address, date of birth, \ldots)
      \item Generate a \textbf{match key (match code)} per record based on
        the criteria
      \item Compare records with identical or similar match codes
      \item Assign a \textbf{confidence score} based on field weights
      \item Records above the threshold are flagged as duplicates
      \item Choose the \textbf{master record}; merge or delete duplicates
    \end{enumerate}
  \end{block}

  \begin{alertblock}{Impact}
    Duplicate customer records lead to: incorrect campaign targeting,
    inflated customer counts, wasted marketing spend.
  \end{alertblock}
\end{frame}

% -------- subsection: Load --------

\begin{frame}
  \frametitle{Load: Bringing Data into the DWH}

  \begin{block}{What the Load Step Does}
    After transformation, cleaned and enriched data is loaded into the
    DWH target tables.
  \end{block}

  \begin{block}{Loading Dimension Tables}
    \begin{itemize}
      \item Look up existing dimension records by natural key
      \item If found: apply SCD logic (Type 1 = update, Type 2 = insert new row)
      \item If not found: insert new record, generate new surrogate key
    \end{itemize}
  \end{block}

  \begin{block}{Loading Fact Tables}
    \begin{itemize}
      \item Resolve all dimension keys (look up surrogate keys)
      \item Append new fact rows (facts are typically never updated)
      \item Handle late-arriving data: facts that arrive after their
        dimension records have been updated
    \end{itemize}
  \end{block}

  \begin{block}{Loading OLAP Cubes}
    \begin{itemize}
      \item Full rebuild of the cube
      \item Or incremental update (add new slices only)
    \end{itemize}
  \end{block}
\end{frame}

% -------- subsection: ETL vs ELT --------

\begin{frame}
  \frametitle{ETL vs.\ ELT}

  \begin{columns}[T]
    \column{0.48\textwidth}
    \begin{block}{Traditional ETL}
      Transform data \textbf{before} loading into the DWH.\\
      Transformation runs in a separate ETL server or tool.\\
      Only clean data enters the DWH.\\
      Typical in on-premise environments.
    \end{block}

    \column{0.48\textwidth}
    \begin{block}{Modern ELT (Extract, Load, Transform)}
      Load raw data \textbf{first}, then transform it \textit{inside}
      the DWH using its compute power.\\
      Raw data is preserved for re-processing.\\
      Enabled by powerful cloud DWH engines (Snowflake, BigQuery).
    \end{block}
  \end{columns}

  \vspace{0.5em}
  \begin{tabularx}{\linewidth}{@{}l X X@{}}
    \toprule
    & \textbf{ETL}                & \textbf{ELT} \\
    \midrule
    Transform where? & External tool         & Inside DWH engine \\
    Data in DWH      & Clean only            & Raw + clean \\
    Flexibility      & Lower                 & Higher \\
    Typical context  & On-premise            & Cloud \\
    Typical tools    & Informatica, SSIS     & dbt, Spark, BigQuery \\
    \bottomrule
  \end{tabularx}
\end{frame}

% -------- subsection: Data Quality Process --------

\begin{frame}
  \frametitle{Data Quality as a Continuous Process}

  Data quality is \textbf{not a one-time fix} --- it requires an ongoing
  improvement cycle.

  % PICTURE PLACEHOLDER
  \begin{alertblock}{Picture: Data Quality Cycle}
    \textit{Add a circular diagram with three phases:
    (1) Analysis (profile and measure quality),
    (2) Model Building (define rules and correction logic),
    (3) Automation (implement rules in ETL pipeline).
    Use circular arrows to show the iterative loop.}
  \end{alertblock}

  \begin{block}{Strategies for Sustained Data Quality}
    \begin{itemize}
      \item Validate at load time; reject or quarantine bad records
      \item Maintain an \textbf{error/quarantine table} for rejected rows
      \item Implement \textbf{data quality dashboards} and metrics
      \item Work with source system owners to fix issues \textit{at the source}
      \item Document \textbf{data lineage}: where does each column come from?
    \end{itemize}
  \end{block}
\end{frame}

% -------- subsection: Metadata --------

\begin{frame}
  \frametitle{Process Modelling and the Role of Metadata}

  \begin{block}{DWH Processes Must Be Managed}
    The ETL pipeline consists of many interdependent processes.
    Managing them requires:
    \begin{itemize}
      \item Definition of \textbf{execution order} and timing
      \item Support for \textbf{parallel processing}
      \item Handling of \textbf{dependencies} between jobs
      \item Definition of \textbf{error handling} and alternative paths
    \end{itemize}
  \end{block}

  \begin{block}{The Role of Metadata in ETL}
    Metadata (especially technical metadata) is the control layer:
    \begin{itemize}
      \item Every process records its \textbf{input and output tables}
      \item ETL jobs can be \textbf{driven by metadata} (rules, mappings stored centrally)
      \item Impact analysis: ``if I change column X, which jobs are affected?''
      \item Change management is much easier with fine-grained, documented processes
    \end{itemize}
  \end{block}
\end{frame}

% -------- subsection: Aggregations --------

\begin{frame}
  \frametitle{Aggregations}

  \begin{block}{What is an Aggregation?}
    An aggregation is the \textbf{summation or consolidation} of data along
    dimension hierarchies.
    For example: summing daily sales to monthly, then annual totals.
  \end{block}

  \begin{block}{Why Pre-Compute Aggregations?}
    \begin{itemize}
      \item Aggregating billions of rows at query time is \textbf{very slow}
      \item Pre-computed aggregations can be served from much smaller tables
      \item The tradeoff: \textbf{storage} for \textbf{speed}
      \item The level of aggregation has a \textbf{huge impact} on data volume
    \end{itemize}
  \end{block}

  \begin{examples}{Simple Aggregation Example}
    \begin{tabularx}{\linewidth}{@{}l l r@{}}
      \toprule
      \textbf{Customer} & \textbf{Product}  & \textbf{Sales (units)} \\
      \midrule
      Meier             & Refrigerator       & 5 \\
      Meier             & Microwave          & 7 \\
      Müller            & Oven               & 6 \\
      Müller            & Refrigerator       & 8 \\
      \midrule
      Meier (total)     &                    & \textbf{12} \\
      Müller (total)    &                    & \textbf{14} \\
      \bottomrule
    \end{tabularx}
  \end{examples}
\end{frame}
